\documentclass[11pt]{article}
\usepackage[T1]{fontenc}
\usepackage{lmodern,microtype}
\usepackage{amsmath,amssymb,booktabs,enumitem}
\usepackage[letterpaper,margin=1in]{geometry}
\usepackage[colorlinks=true,urlcolor=blue,citecolor=blue,linkcolor=blue]{hyperref}
\title{A metaplectic Jacquet--Langlands transfer for division algebras of prime degree, I\\
Revision Log: Draft 8 to Draft 9 (exposition only)}
\author{Prepared by Claude (Anthropic) for David Joyner}
\date{November 8, 2026}
\begin{document}
\maketitle

\section*{Scope}
This is an expository revision requested by David Joyner, the first of
four (one per Part).  \textbf{No mathematical change was made}: no
statement, proof, hypothesis or status label was added, removed or
altered.  A line-by-line \texttt{diff} of Draft~8 and Draft~9 shows
changes only in section and subsection lead-ins, in the pointers of
\S\S1.4--1.6, in the old Sections~10 (Glossary) and~11 (Literature),
which were replaced, and in the bibliography, which was extended.

\paragraph{Numbering is unchanged.}  Parts~II--IV and the companion
paper cite Part~I by number (for example ``Part~I, Lemma~2.2'',
``Proposition~2.3(b)'', ``Proposition~3.5'').  The new material is
therefore placed so that no number moves: the Summary of results is an
unnumbered section between \S1 and \S2, the shared glossary replaces
the old \S10 (Glossary) under the same number, and the full literature
review replaces the old \S11.  We compiled Drafts~8 and~9 and compared
the label tables in their \texttt{.aux} files: every section,
subsection, theorem-like environment, equation and appendix number is
identical.

\section*{Changes}
\begin{enumerate}[leftmargin=*]
\item \textbf{(a) Shared glossary.}  New file \texttt{glossary.tex} at
the top of the repository, written once for the whole series and the
companion paper.  It uses only LaTeX, \texttt{amsmath} and
\texttt{amssymb} (no Part's macros) and labels with the prefix
\texttt{gl:}.  Subsections: status labels (with the rule that labels
are never upgraded without proof and that changes are logged; it
defines \textbf{[P]}, \textbf{[E]}, \textbf{[O]}, \textbf{[C]} with the
wording of Parts~I--IV, and \textbf{[P$^{\mathrm c}$]} with the wording
of the companion paper); global data and the running example; local
fields and the normalization $|p|_p=p^{-1}$; Hilbert symbols; covers
and the cocycle normalization (Kazhdan--Patterson cocycle, parameter
$c$ before and $c_v$ after the relabeling $\zeta\mapsto\zeta^{r'}$, the
ambient cocycle, restriction as an isomorphism of extensions, scalar
centrality, the global normalization by $\kappa_w$ and the splittings,
the star map, and the diagonal commutator matrix of Part~IV, \S2.1);
genuine objects, central characters and good elements; orbital
integrals and matching (including Part~III's notation $J=\Delta O$);
the Haar-measure conventions of Parts~I, II and III, which differ and
are listed side by side; the transform $L_D$ with $b_D$ and $c_D$; and
the named results and hypotheses (Theorems A--E, C$'$, $(\mathrm H_D)$,
(H), $(\mathrm H_D')$, $(\mathrm H')$, (D1)--(D4), Targets A--C).
Every pointer into Parts~II and~III was checked against the label
tables of their current drafts (Part~II, Draft~8: Definitions~2.1 and
6.4, \S6.3; Part~III, Draft~9: Definitions~2.1 and~4.1,
Theorem~4.3).

In Part~I, \S10 is now ``Notation and conventions'': a lead-in,
\verb|% =====================================================================
% glossary.tex
%
% Shared glossary of notation and conventions for the series
%   "A metaplectic Jacquet--Langlands transfer for division algebras
%    of prime degree", Parts I--IV, and the companion paper on the
%    rank-two wild place.
%
% Each Part \input's this file inside its own "Notation and
% conventions" section, for example
%
%   \section{Notation and conventions}\label{sec:glossary}
%   <two-to-four-sentence lead-in>
%   % =====================================================================
% glossary.tex
%
% Shared glossary of notation and conventions for the series
%   "A metaplectic Jacquet--Langlands transfer for division algebras
%    of prime degree", Parts I--IV, and the companion paper on the
%    rank-two wild place.
%
% Each Part \input's this file inside its own "Notation and
% conventions" section, for example
%
%   \section{Notation and conventions}\label{sec:glossary}
%   <two-to-four-sentence lead-in>
%   % =====================================================================
% glossary.tex
%
% Shared glossary of notation and conventions for the series
%   "A metaplectic Jacquet--Langlands transfer for division algebras
%    of prime degree", Parts I--IV, and the companion paper on the
%    rank-two wild place.
%
% Each Part \input's this file inside its own "Notation and
% conventions" section, for example
%
%   \section{Notation and conventions}\label{sec:glossary}
%   <two-to-four-sentence lead-in>
%   \input{../glossary}
%   \subsection{Notation specific to Part I}  ...
%
% The file contains only \subsection's and description lists.  It
% uses only commands from LaTeX, amsmath and amssymb (no macros of
% any Part), so it compiles unchanged in every Part.  Labels carry
% the prefix "gl:" to avoid clashes with the Parts' own labels.
%
% Pointers of the form "Part I, Proposition 4.3" refer to Part I,
% Draft 11 (Part I was reordered in Draft 11; Drafts 8-10 have the
% older numbering, tabulated in Part I's Draft 11 revision log); pointers to Parts II, III
% and IV refer to Drafts 8, 9 and 6.  This file proves nothing.
%
% Version: November 8, 2026 (introduced with Part I, Draft 9; Part I
% pointers renumbered for Part I, Draft 11).
% American English spelling throughout.
% =====================================================================

\subsection{Status labels}\label{gl:labels}
Every statement in the series carries one of the following labels.
A label is never upgraded without a proof, and every change of label
is recorded in the revision log of the draft in which it is made.
\begin{description}
\item[\textbf{[P]}] \emph{Proved here}: a complete argument is
included in the Part in which the label appears, possibly using
results labeled \textbf{[E]}.  A statement proved under an
explicitly stated hypothesis is labeled \textbf{[P]} together with
the hypothesis, as in ``\textbf{[P]} under $(\mathrm H_D)$'', and the
hypothesis itself is listed separately with its own label.
\item[\textbf{[E]}] \emph{Established elsewhere}: a precisely cited
theorem from the literature or from an earlier Part, used as a black
box.  ``\textbf{[P]} (using \textbf{[E]})'' marks a proof that rests
on such inputs; ``\textbf{[E]}/\textbf{[P]}'' marks a statement that
is partly cited and partly proved.
\item[\textbf{[O]}] \emph{Proof outlined}: the steps are identified,
but at least one of them remains to be written out or checked.
\item[\textbf{[C]}] \emph{Conjectural}: new mathematical input is
required.  This includes the goals deferred to later Parts and the
hypotheses under which conditional statements are proved.
\item[\textbf{[P$^{\mathrm c}$]}] \emph{Proved modulo an exhaustive
exact computation}: the proof consists of hand arguments together
with a computation that is exhaustive over an explicitly described
finite set and uses exact arithmetic only (no sampling, and no
floating point without exact re-verification).  Each such script is
committed to the repository, cited by commit hash, and described in
an appendix and in the scripts' README.  At present this label is
used only in the companion paper on the rank-two wild place.
\end{description}

\subsection{Global data}\label{gl:global}
\begin{description}
\item[$K$, $n$, $\mu_n$] $K$ is a totally imaginary number field
containing the group $\mu_n$ of $n$-th roots of unity, and $n\ge1$ is
odd.  All archimedean places of $K$ are complex.
\item[$r$, $D$, $L$, $S$] $r$ is a prime with $(n,r)=1$; $D$ is a
central division algebra over $K$ of dimension $r^2$; $L\subset D$ is
a cyclic maximal subfield, so $[L:K]=r$; $S$ is the (finite) set of
places at which $D$ ramifies.  Since $r$ is prime, every $v\in S$ is
finite and $D_v$ is a division algebra; $|S|\ge2$, and $|S|$ is even
if $r=2$ (Part~I, \S1.1).
\item[$r'$, $c$, $c_v$] Integers with $rr'\equiv1\pmod n$ and
$n\mid r-1+2c$; then $c_v:=r'c$ satisfies $n\mid r-1+2rc_v$ (Part~I,
Assumption~1.1).  Despite the subscript, $c_v$ does not depend on the
place $v$.  The parameter $c$ occurs \emph{before} the relabeling of
$\mu_n$ described in \S\ref{gl:cocycle}, and $c_v$ \emph{after} it.
\item[$\mathbb A_K$, $L_v$] $\mathbb A_K$ is the adele ring of $K$
(written $\mathbb A$ when no confusion is possible), and
$L_v:=L\otimes_KK_v=\prod_{w\mid v}L_w$.
\item[$G_D$] The algebraic group $D^\times$ over $K$, so that
$G_D(K_v)=D_v^\times$.
\item[Running example] $K=\mathbb Q(\zeta_3)$, $n=3$, $r=2$,
$D=(-1,4-\zeta_3)_K$, ramified exactly at
$S=\{\mathfrak p_3,\mathfrak p_7\}$ with $\mathfrak p_3=(1-\zeta_3)$
and $\mathfrak p_7=(3+\zeta_3)$; $L=K(i)=\mathbb Q(\zeta_{12})$,
$c=1$, $r'=2$, $c_v=2$.  The place $\mathfrak p_3$ divides $n$ and
$\mathfrak p_7$ does not.  The note
\texttt{notes/mjl\_running\_example.tex} works the example out across
the series.
\end{description}

\subsection{Local fields and absolute values}\label{gl:local}
\begin{description}
\item[$F$, $\mathcal O_F$, $\varpi$, $q_F$, $p$] $F$ is a
nonarchimedean local field of characteristic zero containing $\mu_n$
(usually $F=K_v$), with ring of integers $\mathcal O_F$, uniformizer
$\varpi$, residue field of order $q_F$, and residual characteristic
$p$.
\item[$|\cdot|_F$, $|\cdot|_p$] $|\cdot|_F$ is the normalized absolute
value, $|\varpi|_F=q_F^{-1}$; on $\mathbb Q_p$ it is written
$|\cdot|_p$, with $|p|_p=p^{-1}$.  In particular $|n|_F=1$ if
$p\nmid n$, and $|F^\times/F^{\times n}|=n^2|n|_F^{-1}$ (Part~I,
proof of Proposition~5.5).
\item[Places dividing $n$; tame and wild places] A place $v$ (or the
field $F=K_v$) \emph{divides $n$} if $p\mid n$.  The places $v\in S$
with $p\nmid n$ are called \emph{tame}, and those with $p\mid n$
\emph{wild}.  In the running example $\mathfrak p_7$ is tame and
$\mathfrak p_3$ is wild ($|3|_{\mathfrak p_3}=3^{-2}$).
\end{description}

\subsection{Hilbert symbols}\label{gl:hilbert}
\begin{description}
\item[$(a,b)_F$] The $n$-th Hilbert symbol of $F$, with values in
$\mu_n$; Part~IV and the companion paper also write $(a,b)_{n,F}$ or
$(a,b)_n$.  It is bimultiplicative, nondegenerate and antisymmetric,
$(a,b)_F=(b,a)_F^{-1}$.  At complex places it is trivial.
\item[Norm compatibility] For a finite extension $F'/F$, $a\in
F^\times$ and $b\in F'^\times$, $(a,b)_{F'}=(a,N_{F'/F}b)_F$; in
particular $(a,b)_{F'}=(a,b)_F^{[F':F]}$ for $a,b\in F^\times$
[Serre, \emph{Local Fields}, Ch.~XIV].
\item[$(a,b)_{\mathbb A}$] For a number field $M\supseteq K$,
$(a,b)_{\mathbb A_M}=\prod_w(a_w,b_w)_{M_w}$.  Hilbert reciprocity:
$(a,b)_{\mathbb A_M}=1$ for $a,b\in M^\times$.
\end{description}

\subsection{Covers and the cocycle normalization}\label{gl:cocycle}
\begin{description}
\item[Cover, $p$, $\overline{H'}$] An \emph{$n$-fold cover} of a
locally compact group $H$ is a topological central extension
$1\to\mu_n\to\overline H\xrightarrow{\ p\ }H\to1$.  For a subgroup
$H'\subset H$ we write $\overline{H'}:=p^{-1}(H')$.  Parts~I and~II
write $\overline H$; Parts~III, IV and the companion paper write
$\widetilde H$, $\widetilde Z$.
\item[Commutators] For lifts $\tilde x,\tilde y$ of commuting
elements, $[x,y]:=\tilde x\tilde y\tilde x^{-1}\tilde y^{-1}\in\mu_n$;
it does not depend on the lifts, and it is bimultiplicative on a
torus [Flicker--Kazhdan, \S3].
\item[Kazhdan--Patterson cocycle $\sigma_F$] The cocycle of
[Kazhdan--Patterson, \S0.1] and [Flicker--Kazhdan, \S2] on
$\mathrm{GL}(r,F)$ with twisting parameter $0$.  On the diagonal
torus,
$\sigma_F(\mathrm{diag}(a_i),\mathrm{diag}(b_i))=\prod_{i<j}(a_i,b_j)_F$.
\item[Parameter-$c'$ cover $\overline G^{(c')}$] The cover of
$\mathrm{GL}(r,F)$ with cocycle
$\sigma_F(g,h)\,(\det g,\det h)_F^{\,c'}$.  The comparison uses
$c'=c_v$; the parameter is forced (Part~I, Remark~5.3).  On the
diagonal torus the commutator of $\mathrm{diag}(t_i)$ and
$\mathrm{diag}(u_j)$ is $\prod_{i,j}(t_i,u_j)_F^{A_{ij}}$ with
$A=-I_r+(1+2c_v)J_r$, where $J_r$ is the all-ones matrix (Part~IV,
\S2.1).
\item[Ambient cover $\overline{\mathcal G}_v$] The cover of
$\mathcal G_v=\mathrm{GL}(r,L_v)=\prod_{w\mid v}\mathrm{GL}(r,L_w)$
with cocycle
$\prod_{w\mid v}\sigma_{L_w}(g_w,h_w)\,(\det g_w,\det h_w)_{L_w}^{\,c_v}$,
the copies of $\mu_n$ being identified (Part~I, \S4.2 and
equation~(1)).  Both covers compared in the series are subgroups of
it, so their common central subgroup $p^{-1}(K_v^\times)$ is literally
the same subgroup on the two sides.
\item[Restriction and relabeling] Pulled back along
$\iota_D:D_v^\times\hookrightarrow\mathcal G_v$, the ambient cocycle
is the local component of the cocycle
$\beta=\sigma_L\cdot(\mathrm{Nrd},\mathrm{Nrd})^c$ of the corrected
1988 paper [JC, equation~(4)].  Pulled back along the diagonal
embedding of $\mathrm{GL}(r,K_v)$, it is isomorphic, \emph{as an
extension}, to the extension with cocycle
$\sigma_F^{\,r}(\det,\det)_F^{\,c}$ (Part~I, Proposition~4.3).  The
group $\mu_n$ is relabeled \emph{once}, at every place and globally,
by $\zeta\mapsto\zeta^{r'}$; after relabeling the second extension is
$\overline G^{(c_v)}$.  Formulas state whether they hold before or
after relabeling.  The isomorphism is one of extensions, not an
equality of cocycles, and it is unique up to
$\tilde g\mapsto\chi(\det p(\tilde g))\,\tilde g$ with
$\chi:F^\times\to\mu_n$ continuous.  Central characters are therefore
never transported through it (Part~I, Remark~4.5).
\item[Scalar centrality] Because $n\mid r-1+2rc_v$, the whole inverse
image $\overline Z=p^{-1}(F^\times)$ of the scalars is central
(Part~I, Lemma~4.2).  A cover with this property is called
\emph{scalar-central}.
\item[Global normalization] Globally one uses the normalized cocycle
$\tau_w(g,h)=\sigma_w(g,h)\,\kappa_w(g)\kappa_w(h)\kappa_w(gh)^{-1}$
of [Flicker--Kazhdan, \S2], where $\kappa_w$ is trivial for $w\mid
n\infty$.  The compact splitting is $\kappa^*_w(k)=(k,1)_\tau$ on
$\mathrm{GL}(r,\mathcal O_{L_w})$ for $w\nmid n\infty$, and the
rational splitting is $\iota(x)=s_\tau(x)\,\kappa_L(x)^{-1}$; in the
unnormalized models the rational splitting is $x\mapsto(x,1)_\sigma$
at every place (Part~I, Lemma~7.1).  The global ambient cover
$\overline{\mathcal G}(\mathbb A_L)$ has cocycle
$\tau_L\cdot(\det,\det)_{\mathbb A_L}^{\,c_v}$.
\item[Star map] For $h$ in the base group and any lift $\tilde h$,
$h^*:=\tilde h^{\,n}$.  It does not depend on the lift, it is
conjugation equivariant, $(ah)^*=a^*h^*$ for scalars $a$, and it is
multiplicative on every torus because $n$ is odd (Part~III,
Definition~2.1; Part~IV, \S2.2).  For $\zeta\in\mu_n(F)$ one has
$\zeta^*=1$.
\end{description}

\subsection{Genuine objects, central characters and good elements}\label{gl:genuine}
\begin{description}
\item[Genuine representation] A representation of $\overline H$ on
which $\mu_n$ acts by the tautological character $\zeta\mapsto\zeta$.
\item[$\tilde\theta$, $\theta$] A genuine character $\tilde\theta$ of
$\overline Z$ (unitary from Part~I, \S7.1 on; globally, a character of
$\overline Z(\mathbb A_K)$ trivial on $\iota(K^\times)$).  The same
$\tilde\theta$ serves both covers, since $\overline Z$ is a common
subgroup.  The ordinary character attached to it is
$\theta(a):=\tilde\theta(a^*)$; it is trivial on $\mu_n$, and
$\theta$ does not determine $\tilde\theta$ (Part~IV, \S2.2).
Part~III requires only that $\tilde\theta|_{\mu_n}$ be faithful.
\item[Genuine test functions] $C_c^\infty(\overline H,\tilde\theta)_{\mathrm{gen}}$
is the space of smooth functions with support compact modulo
$\overline Z$ and $f(z\tilde h)=\tilde\theta(z)^{-1}f(\tilde h)$ for
$z\in\overline Z$; in particular $f(\zeta\tilde h)=\zeta^{-1}f(\tilde h)$.
Characters and invariant class functions of genuine representations
transform by $\tilde\theta$ itself (Part~III, \S2).
\item[Regular, elliptic] An element of $\mathrm{GL}(r,F)$ is
\emph{regular} if its eigenvalues are distinct, and \emph{elliptic
regular} if moreover its characteristic polynomial is irreducible.
Since $r$ is prime, the regular elements of $D_v^\times$ are the
noncentral ones, and they are all elliptic.  An element of a cover is
regular (elliptic) if its image is.
\item[Good element] A regular $\tilde\gamma$ is \emph{good} if it
commutes with $p^{-1}(T_\gamma)$, where $T_\gamma$ is the centralizer
of $\gamma=p(\tilde\gamma)$; equivalently, $\tilde\gamma$ is central
in the covered torus $\overline{T_\gamma}$.  Genuine invariant
functions vanish at regular elements that are not good.  For a
maximal subfield $E$ of $D_v$ (or an elliptic torus $E^\times$ of
$\mathrm{GL}(r,F)$), the good elements of $\overline{E^\times}$ are
$p^{-1}(F^\times E^{\times n})=\overline Z\cdot\{t^*\}$, of index
$|n|_F^{-(r-1)}$, and every genuine irreducible representation of
$\overline{E^\times}$ has dimension $|n|_F^{-(r-1)/2}$ (Part~I,
Proposition~5.5).  Covered elliptic tori are abelian exactly at the
places not dividing $n$.
\end{description}

\subsection{Orbital integrals and matching}\label{gl:matching}
\begin{description}
\item[$\Delta$] For regular semisimple $\gamma$,
$\Delta(\gamma)=|D(\gamma)|_F^{1/2}|\det\gamma|_F^{-(r-1)/2}$, where
$D(\gamma)$ is the discriminant of the (reduced) characteristic
polynomial and $\det=\mathrm{Nrd}$ on the division-algebra side.
Thus $\Delta(\gamma)^2$ is the Weyl discriminant
$|\det(1-\mathrm{Ad}(\gamma^{-1}))|_{\mathfrak g/\mathfrak t}|_F$
(Part~I, Definition~6.1 and Lemma~6.12).
\item[$\Phi$, $F=\Delta\Phi$] $\Phi(\tilde\gamma,\tilde f)$ is the
orbital integral over $T_\gamma\backslash G$ of a genuine test
function at a good regular $\tilde\gamma$, and
$F(\tilde\gamma,\tilde f)=\Delta(\gamma)\Phi(\tilde\gamma,\tilde f)$
is the normalized orbital integral, set to $0$ at regular elements
that are not good (Part~I, Definition~6.1).  Part~III writes
$J(t,f)=\Delta(t)\,O(t,f)$ for the same function, with its own
measure conventions.
\item[Canonical identification $\jmath$] At $v\in S$ and for a maximal
subfield $E\subset D_v$ with an embedding $j:E\hookrightarrow
M_r(F)$, every isomorphism of covered tori that is the identity on
$\overline Z$ has the same restriction $\jmath_{E,j}$ to the good
elements; these restrictions induce a bijection $\jmath$ from the good
regular classes of the cover of $D_v^\times$ to the good regular
elliptic classes of the cover of $\mathrm{GL}(r,F)$ (Part~I,
Theorem~5.9 and Corollary~5.12).
\item[Matching] Genuine test functions $\tilde f_D$ and $\tilde f$
\emph{match} if (M1)
$F(\tilde\gamma,\tilde f_D)=(-1)^{r-1}F(\jmath(\tilde\gamma),\tilde f)$
at every good regular $\tilde\gamma$, and (M2) $F(\tilde\delta,\tilde
f)=0$ at every good regular non-elliptic $\tilde\delta$ (Part~I,
Definition~6.3).  The sign sits in the matching of functions, so that
the expected character relation is
$\Theta_{\overline\pi_D}=(-1)^{r-1}\Theta_{\overline\pi}\circ\jmath$.
\item[Supercusp form] A genuine test function on the cover of
$\mathrm{GL}(r,F)$ whose integrals along the unipotent radicals of
all proper parabolic subgroups vanish, for example a matrix
coefficient of a supercuspidal representation [Flicker--Kazhdan,
\S13].
\item[Germ] The asymptotic behavior of orbital integrals near a
singular element, expressed through Shalika germs [Flicker--Kazhdan,
\S8]; used near the center in Part~II.
\end{description}

\subsection{Haar measures}\label{gl:measures}
\begin{description}
\item[Part~I] At $v\in S$ the Haar measures on $D_v^\times/F^\times$
and $\mathrm{GL}(r,F)/F^\times$ are arbitrary but fixed, and matching
is defined with respect to them; torus measures are transported by
conjugation and do not affect matching.  At $v\notin S$ the measures
correspond under the identifications $\psi_v$ (Part~I, \S\S6.1, 7.3).
\item[Part~II] At $v\in S$,
$\mathrm{vol}(D_v^\times/F^\times)=d(\mathrm{St})^{-1}$, where
$d(\mathrm{St})$ is the formal degree of the Steinberg representation
for the chosen measure on $\mathrm{GL}(r,F)/F^\times$ (Part~II,
Definition~2.1).
\item[Part~III] For the transform $L_D$, the quotients $H/Z$ and
$T/Z$ carry their Haar \emph{probability} measures (Part~III, \S4).
\item[Formal degree] For a representation $\sigma$ that is square
integrable modulo the center, the constant $d(\sigma)$ in the Schur
orthogonality relations; it depends on the Haar measure.
\end{description}

\subsection{The division-side transform}\label{gl:transform}
\begin{description}
\item[$H$, $\widetilde H$, $L_D$] For a local division algebra $D$ of
degree $r$ over $F$, $H=D^\times$ and $\widetilde H$ is its cover from
Part~I.  The transform $L_D$ sends $\theta$-equivariant invariant
$L^2$ functions $\varphi$ on $H$ to genuine
$\tilde\theta$-equivariant invariant functions on $\widetilde H$ by
$(L_D\varphi)(z\,t^*)=b_D\,\tilde\theta(z)\,\Delta(t)\Delta(t^n)^{-1}\varphi(t)$
at good regular elements, and $0$ at the others; it is unitary
(Part~III, Definition~4.1 and Theorem~4.3).
\item[$b_D$, $c_D$] $b_D=|n|_F^{-(r-1)/2}$ and
$c_D=|n|_F^{-(r^2-1)/2}$ (Part~III, \S\S4--5).  Both equal $1$ at a
tame place; in the running example $b_D=3$ and $c_D=27$ at
$\mathfrak p_3$.
\item[Integrality of $L_D$] The question whether $L_D$ carries
irreducible characters to irreducible genuine characters.  It holds at
tame places (Part~III) and is open in general at wild places
(Part~IV; companion paper).
\end{description}

\subsection{Named results and hypotheses}\label{gl:named}
\begin{description}
\item[Theorems A, B, C] Part~I: the comparison of covered elliptic
tori, the matching of regular-supported test functions, and the
global transfer under $(\mathrm H_D)$ (with its converse under (H)).
\item[Theorems D, E, C$'$] Part~II: the equivariant section near the
center and its compatibility with $\jmath$; the matching of arbitrary
genuine test functions; and the global transfer under the weaker
hypothesis $(\mathrm H_D')$ (with its converse under $(\mathrm H')$).
\item[$(\mathrm H_D)$, (H), $(\mathrm H_D')$, $(\mathrm H')$]
Non-cancellation hypotheses: the character of a near-equivalence
packet at the places of $S$ is not identically zero on the relevant
matching test functions (Part~I, Definition~9.9 and
Theorem~9.12; Part~II, Definition~6.4 and \S6.3).  $(\mathrm H_D)$
holds when the $D$-side packet is finite (Part~I,
Proposition~9.10); sufficient conditions for (H) are in Part~I,
Remark~9.13; and $(\mathrm H_D)\Rightarrow(\mathrm H_D')$,
$(\mathrm H)\Rightarrow(\mathrm H')$ (Part~II, \S6.3).
\item[(D1)--(D4)] The steps left open by Part~I: matching of
arbitrary functions, local character relations, finiteness of
near-equivalence packets, and composition with [JC] (Part~I, \S10).
Their status in the later Parts is recorded in Part~I, \S1.3.
\item[Targets A, B, C] The companion paper's three local problems at
the wild place of the running example and beyond (companion paper,
\S1).
\end{description}

%   \subsection{Notation specific to Part I}  ...
%
% The file contains only \subsection's and description lists.  It
% uses only commands from LaTeX, amsmath and amssymb (no macros of
% any Part), so it compiles unchanged in every Part.  Labels carry
% the prefix "gl:" to avoid clashes with the Parts' own labels.
%
% Pointers of the form "Part I, Proposition 4.3" refer to Part I,
% Draft 11 (Part I was reordered in Draft 11; Drafts 8-10 have the
% older numbering, tabulated in Part I's Draft 11 revision log); pointers to Parts II, III
% and IV refer to Drafts 8, 9 and 6.  This file proves nothing.
%
% Version: November 8, 2026 (introduced with Part I, Draft 9; Part I
% pointers renumbered for Part I, Draft 11).
% American English spelling throughout.
% =====================================================================

\subsection{Status labels}\label{gl:labels}
Every statement in the series carries one of the following labels.
A label is never upgraded without a proof, and every change of label
is recorded in the revision log of the draft in which it is made.
\begin{description}
\item[\textbf{[P]}] \emph{Proved here}: a complete argument is
included in the Part in which the label appears, possibly using
results labeled \textbf{[E]}.  A statement proved under an
explicitly stated hypothesis is labeled \textbf{[P]} together with
the hypothesis, as in ``\textbf{[P]} under $(\mathrm H_D)$'', and the
hypothesis itself is listed separately with its own label.
\item[\textbf{[E]}] \emph{Established elsewhere}: a precisely cited
theorem from the literature or from an earlier Part, used as a black
box.  ``\textbf{[P]} (using \textbf{[E]})'' marks a proof that rests
on such inputs; ``\textbf{[E]}/\textbf{[P]}'' marks a statement that
is partly cited and partly proved.
\item[\textbf{[O]}] \emph{Proof outlined}: the steps are identified,
but at least one of them remains to be written out or checked.
\item[\textbf{[C]}] \emph{Conjectural}: new mathematical input is
required.  This includes the goals deferred to later Parts and the
hypotheses under which conditional statements are proved.
\item[\textbf{[P$^{\mathrm c}$]}] \emph{Proved modulo an exhaustive
exact computation}: the proof consists of hand arguments together
with a computation that is exhaustive over an explicitly described
finite set and uses exact arithmetic only (no sampling, and no
floating point without exact re-verification).  Each such script is
committed to the repository, cited by commit hash, and described in
an appendix and in the scripts' README.  At present this label is
used only in the companion paper on the rank-two wild place.
\end{description}

\subsection{Global data}\label{gl:global}
\begin{description}
\item[$K$, $n$, $\mu_n$] $K$ is a totally imaginary number field
containing the group $\mu_n$ of $n$-th roots of unity, and $n\ge1$ is
odd.  All archimedean places of $K$ are complex.
\item[$r$, $D$, $L$, $S$] $r$ is a prime with $(n,r)=1$; $D$ is a
central division algebra over $K$ of dimension $r^2$; $L\subset D$ is
a cyclic maximal subfield, so $[L:K]=r$; $S$ is the (finite) set of
places at which $D$ ramifies.  Since $r$ is prime, every $v\in S$ is
finite and $D_v$ is a division algebra; $|S|\ge2$, and $|S|$ is even
if $r=2$ (Part~I, \S1.1).
\item[$r'$, $c$, $c_v$] Integers with $rr'\equiv1\pmod n$ and
$n\mid r-1+2c$; then $c_v:=r'c$ satisfies $n\mid r-1+2rc_v$ (Part~I,
Assumption~1.1).  Despite the subscript, $c_v$ does not depend on the
place $v$.  The parameter $c$ occurs \emph{before} the relabeling of
$\mu_n$ described in \S\ref{gl:cocycle}, and $c_v$ \emph{after} it.
\item[$\mathbb A_K$, $L_v$] $\mathbb A_K$ is the adele ring of $K$
(written $\mathbb A$ when no confusion is possible), and
$L_v:=L\otimes_KK_v=\prod_{w\mid v}L_w$.
\item[$G_D$] The algebraic group $D^\times$ over $K$, so that
$G_D(K_v)=D_v^\times$.
\item[Running example] $K=\mathbb Q(\zeta_3)$, $n=3$, $r=2$,
$D=(-1,4-\zeta_3)_K$, ramified exactly at
$S=\{\mathfrak p_3,\mathfrak p_7\}$ with $\mathfrak p_3=(1-\zeta_3)$
and $\mathfrak p_7=(3+\zeta_3)$; $L=K(i)=\mathbb Q(\zeta_{12})$,
$c=1$, $r'=2$, $c_v=2$.  The place $\mathfrak p_3$ divides $n$ and
$\mathfrak p_7$ does not.  The note
\texttt{notes/mjl\_running\_example.tex} works the example out across
the series.
\end{description}

\subsection{Local fields and absolute values}\label{gl:local}
\begin{description}
\item[$F$, $\mathcal O_F$, $\varpi$, $q_F$, $p$] $F$ is a
nonarchimedean local field of characteristic zero containing $\mu_n$
(usually $F=K_v$), with ring of integers $\mathcal O_F$, uniformizer
$\varpi$, residue field of order $q_F$, and residual characteristic
$p$.
\item[$|\cdot|_F$, $|\cdot|_p$] $|\cdot|_F$ is the normalized absolute
value, $|\varpi|_F=q_F^{-1}$; on $\mathbb Q_p$ it is written
$|\cdot|_p$, with $|p|_p=p^{-1}$.  In particular $|n|_F=1$ if
$p\nmid n$, and $|F^\times/F^{\times n}|=n^2|n|_F^{-1}$ (Part~I,
proof of Proposition~5.5).
\item[Places dividing $n$; tame and wild places] A place $v$ (or the
field $F=K_v$) \emph{divides $n$} if $p\mid n$.  The places $v\in S$
with $p\nmid n$ are called \emph{tame}, and those with $p\mid n$
\emph{wild}.  In the running example $\mathfrak p_7$ is tame and
$\mathfrak p_3$ is wild ($|3|_{\mathfrak p_3}=3^{-2}$).
\end{description}

\subsection{Hilbert symbols}\label{gl:hilbert}
\begin{description}
\item[$(a,b)_F$] The $n$-th Hilbert symbol of $F$, with values in
$\mu_n$; Part~IV and the companion paper also write $(a,b)_{n,F}$ or
$(a,b)_n$.  It is bimultiplicative, nondegenerate and antisymmetric,
$(a,b)_F=(b,a)_F^{-1}$.  At complex places it is trivial.
\item[Norm compatibility] For a finite extension $F'/F$, $a\in
F^\times$ and $b\in F'^\times$, $(a,b)_{F'}=(a,N_{F'/F}b)_F$; in
particular $(a,b)_{F'}=(a,b)_F^{[F':F]}$ for $a,b\in F^\times$
[Serre, \emph{Local Fields}, Ch.~XIV].
\item[$(a,b)_{\mathbb A}$] For a number field $M\supseteq K$,
$(a,b)_{\mathbb A_M}=\prod_w(a_w,b_w)_{M_w}$.  Hilbert reciprocity:
$(a,b)_{\mathbb A_M}=1$ for $a,b\in M^\times$.
\end{description}

\subsection{Covers and the cocycle normalization}\label{gl:cocycle}
\begin{description}
\item[Cover, $p$, $\overline{H'}$] An \emph{$n$-fold cover} of a
locally compact group $H$ is a topological central extension
$1\to\mu_n\to\overline H\xrightarrow{\ p\ }H\to1$.  For a subgroup
$H'\subset H$ we write $\overline{H'}:=p^{-1}(H')$.  Parts~I and~II
write $\overline H$; Parts~III, IV and the companion paper write
$\widetilde H$, $\widetilde Z$.
\item[Commutators] For lifts $\tilde x,\tilde y$ of commuting
elements, $[x,y]:=\tilde x\tilde y\tilde x^{-1}\tilde y^{-1}\in\mu_n$;
it does not depend on the lifts, and it is bimultiplicative on a
torus [Flicker--Kazhdan, \S3].
\item[Kazhdan--Patterson cocycle $\sigma_F$] The cocycle of
[Kazhdan--Patterson, \S0.1] and [Flicker--Kazhdan, \S2] on
$\mathrm{GL}(r,F)$ with twisting parameter $0$.  On the diagonal
torus,
$\sigma_F(\mathrm{diag}(a_i),\mathrm{diag}(b_i))=\prod_{i<j}(a_i,b_j)_F$.
\item[Parameter-$c'$ cover $\overline G^{(c')}$] The cover of
$\mathrm{GL}(r,F)$ with cocycle
$\sigma_F(g,h)\,(\det g,\det h)_F^{\,c'}$.  The comparison uses
$c'=c_v$; the parameter is forced (Part~I, Remark~5.3).  On the
diagonal torus the commutator of $\mathrm{diag}(t_i)$ and
$\mathrm{diag}(u_j)$ is $\prod_{i,j}(t_i,u_j)_F^{A_{ij}}$ with
$A=-I_r+(1+2c_v)J_r$, where $J_r$ is the all-ones matrix (Part~IV,
\S2.1).
\item[Ambient cover $\overline{\mathcal G}_v$] The cover of
$\mathcal G_v=\mathrm{GL}(r,L_v)=\prod_{w\mid v}\mathrm{GL}(r,L_w)$
with cocycle
$\prod_{w\mid v}\sigma_{L_w}(g_w,h_w)\,(\det g_w,\det h_w)_{L_w}^{\,c_v}$,
the copies of $\mu_n$ being identified (Part~I, \S4.2 and
equation~(1)).  Both covers compared in the series are subgroups of
it, so their common central subgroup $p^{-1}(K_v^\times)$ is literally
the same subgroup on the two sides.
\item[Restriction and relabeling] Pulled back along
$\iota_D:D_v^\times\hookrightarrow\mathcal G_v$, the ambient cocycle
is the local component of the cocycle
$\beta=\sigma_L\cdot(\mathrm{Nrd},\mathrm{Nrd})^c$ of the corrected
1988 paper [JC, equation~(4)].  Pulled back along the diagonal
embedding of $\mathrm{GL}(r,K_v)$, it is isomorphic, \emph{as an
extension}, to the extension with cocycle
$\sigma_F^{\,r}(\det,\det)_F^{\,c}$ (Part~I, Proposition~4.3).  The
group $\mu_n$ is relabeled \emph{once}, at every place and globally,
by $\zeta\mapsto\zeta^{r'}$; after relabeling the second extension is
$\overline G^{(c_v)}$.  Formulas state whether they hold before or
after relabeling.  The isomorphism is one of extensions, not an
equality of cocycles, and it is unique up to
$\tilde g\mapsto\chi(\det p(\tilde g))\,\tilde g$ with
$\chi:F^\times\to\mu_n$ continuous.  Central characters are therefore
never transported through it (Part~I, Remark~4.5).
\item[Scalar centrality] Because $n\mid r-1+2rc_v$, the whole inverse
image $\overline Z=p^{-1}(F^\times)$ of the scalars is central
(Part~I, Lemma~4.2).  A cover with this property is called
\emph{scalar-central}.
\item[Global normalization] Globally one uses the normalized cocycle
$\tau_w(g,h)=\sigma_w(g,h)\,\kappa_w(g)\kappa_w(h)\kappa_w(gh)^{-1}$
of [Flicker--Kazhdan, \S2], where $\kappa_w$ is trivial for $w\mid
n\infty$.  The compact splitting is $\kappa^*_w(k)=(k,1)_\tau$ on
$\mathrm{GL}(r,\mathcal O_{L_w})$ for $w\nmid n\infty$, and the
rational splitting is $\iota(x)=s_\tau(x)\,\kappa_L(x)^{-1}$; in the
unnormalized models the rational splitting is $x\mapsto(x,1)_\sigma$
at every place (Part~I, Lemma~7.1).  The global ambient cover
$\overline{\mathcal G}(\mathbb A_L)$ has cocycle
$\tau_L\cdot(\det,\det)_{\mathbb A_L}^{\,c_v}$.
\item[Star map] For $h$ in the base group and any lift $\tilde h$,
$h^*:=\tilde h^{\,n}$.  It does not depend on the lift, it is
conjugation equivariant, $(ah)^*=a^*h^*$ for scalars $a$, and it is
multiplicative on every torus because $n$ is odd (Part~III,
Definition~2.1; Part~IV, \S2.2).  For $\zeta\in\mu_n(F)$ one has
$\zeta^*=1$.
\end{description}

\subsection{Genuine objects, central characters and good elements}\label{gl:genuine}
\begin{description}
\item[Genuine representation] A representation of $\overline H$ on
which $\mu_n$ acts by the tautological character $\zeta\mapsto\zeta$.
\item[$\tilde\theta$, $\theta$] A genuine character $\tilde\theta$ of
$\overline Z$ (unitary from Part~I, \S7.1 on; globally, a character of
$\overline Z(\mathbb A_K)$ trivial on $\iota(K^\times)$).  The same
$\tilde\theta$ serves both covers, since $\overline Z$ is a common
subgroup.  The ordinary character attached to it is
$\theta(a):=\tilde\theta(a^*)$; it is trivial on $\mu_n$, and
$\theta$ does not determine $\tilde\theta$ (Part~IV, \S2.2).
Part~III requires only that $\tilde\theta|_{\mu_n}$ be faithful.
\item[Genuine test functions] $C_c^\infty(\overline H,\tilde\theta)_{\mathrm{gen}}$
is the space of smooth functions with support compact modulo
$\overline Z$ and $f(z\tilde h)=\tilde\theta(z)^{-1}f(\tilde h)$ for
$z\in\overline Z$; in particular $f(\zeta\tilde h)=\zeta^{-1}f(\tilde h)$.
Characters and invariant class functions of genuine representations
transform by $\tilde\theta$ itself (Part~III, \S2).
\item[Regular, elliptic] An element of $\mathrm{GL}(r,F)$ is
\emph{regular} if its eigenvalues are distinct, and \emph{elliptic
regular} if moreover its characteristic polynomial is irreducible.
Since $r$ is prime, the regular elements of $D_v^\times$ are the
noncentral ones, and they are all elliptic.  An element of a cover is
regular (elliptic) if its image is.
\item[Good element] A regular $\tilde\gamma$ is \emph{good} if it
commutes with $p^{-1}(T_\gamma)$, where $T_\gamma$ is the centralizer
of $\gamma=p(\tilde\gamma)$; equivalently, $\tilde\gamma$ is central
in the covered torus $\overline{T_\gamma}$.  Genuine invariant
functions vanish at regular elements that are not good.  For a
maximal subfield $E$ of $D_v$ (or an elliptic torus $E^\times$ of
$\mathrm{GL}(r,F)$), the good elements of $\overline{E^\times}$ are
$p^{-1}(F^\times E^{\times n})=\overline Z\cdot\{t^*\}$, of index
$|n|_F^{-(r-1)}$, and every genuine irreducible representation of
$\overline{E^\times}$ has dimension $|n|_F^{-(r-1)/2}$ (Part~I,
Proposition~5.5).  Covered elliptic tori are abelian exactly at the
places not dividing $n$.
\end{description}

\subsection{Orbital integrals and matching}\label{gl:matching}
\begin{description}
\item[$\Delta$] For regular semisimple $\gamma$,
$\Delta(\gamma)=|D(\gamma)|_F^{1/2}|\det\gamma|_F^{-(r-1)/2}$, where
$D(\gamma)$ is the discriminant of the (reduced) characteristic
polynomial and $\det=\mathrm{Nrd}$ on the division-algebra side.
Thus $\Delta(\gamma)^2$ is the Weyl discriminant
$|\det(1-\mathrm{Ad}(\gamma^{-1}))|_{\mathfrak g/\mathfrak t}|_F$
(Part~I, Definition~6.1 and Lemma~6.12).
\item[$\Phi$, $F=\Delta\Phi$] $\Phi(\tilde\gamma,\tilde f)$ is the
orbital integral over $T_\gamma\backslash G$ of a genuine test
function at a good regular $\tilde\gamma$, and
$F(\tilde\gamma,\tilde f)=\Delta(\gamma)\Phi(\tilde\gamma,\tilde f)$
is the normalized orbital integral, set to $0$ at regular elements
that are not good (Part~I, Definition~6.1).  Part~III writes
$J(t,f)=\Delta(t)\,O(t,f)$ for the same function, with its own
measure conventions.
\item[Canonical identification $\jmath$] At $v\in S$ and for a maximal
subfield $E\subset D_v$ with an embedding $j:E\hookrightarrow
M_r(F)$, every isomorphism of covered tori that is the identity on
$\overline Z$ has the same restriction $\jmath_{E,j}$ to the good
elements; these restrictions induce a bijection $\jmath$ from the good
regular classes of the cover of $D_v^\times$ to the good regular
elliptic classes of the cover of $\mathrm{GL}(r,F)$ (Part~I,
Theorem~5.9 and Corollary~5.12).
\item[Matching] Genuine test functions $\tilde f_D$ and $\tilde f$
\emph{match} if (M1)
$F(\tilde\gamma,\tilde f_D)=(-1)^{r-1}F(\jmath(\tilde\gamma),\tilde f)$
at every good regular $\tilde\gamma$, and (M2) $F(\tilde\delta,\tilde
f)=0$ at every good regular non-elliptic $\tilde\delta$ (Part~I,
Definition~6.3).  The sign sits in the matching of functions, so that
the expected character relation is
$\Theta_{\overline\pi_D}=(-1)^{r-1}\Theta_{\overline\pi}\circ\jmath$.
\item[Supercusp form] A genuine test function on the cover of
$\mathrm{GL}(r,F)$ whose integrals along the unipotent radicals of
all proper parabolic subgroups vanish, for example a matrix
coefficient of a supercuspidal representation [Flicker--Kazhdan,
\S13].
\item[Germ] The asymptotic behavior of orbital integrals near a
singular element, expressed through Shalika germs [Flicker--Kazhdan,
\S8]; used near the center in Part~II.
\end{description}

\subsection{Haar measures}\label{gl:measures}
\begin{description}
\item[Part~I] At $v\in S$ the Haar measures on $D_v^\times/F^\times$
and $\mathrm{GL}(r,F)/F^\times$ are arbitrary but fixed, and matching
is defined with respect to them; torus measures are transported by
conjugation and do not affect matching.  At $v\notin S$ the measures
correspond under the identifications $\psi_v$ (Part~I, \S\S6.1, 7.3).
\item[Part~II] At $v\in S$,
$\mathrm{vol}(D_v^\times/F^\times)=d(\mathrm{St})^{-1}$, where
$d(\mathrm{St})$ is the formal degree of the Steinberg representation
for the chosen measure on $\mathrm{GL}(r,F)/F^\times$ (Part~II,
Definition~2.1).
\item[Part~III] For the transform $L_D$, the quotients $H/Z$ and
$T/Z$ carry their Haar \emph{probability} measures (Part~III, \S4).
\item[Formal degree] For a representation $\sigma$ that is square
integrable modulo the center, the constant $d(\sigma)$ in the Schur
orthogonality relations; it depends on the Haar measure.
\end{description}

\subsection{The division-side transform}\label{gl:transform}
\begin{description}
\item[$H$, $\widetilde H$, $L_D$] For a local division algebra $D$ of
degree $r$ over $F$, $H=D^\times$ and $\widetilde H$ is its cover from
Part~I.  The transform $L_D$ sends $\theta$-equivariant invariant
$L^2$ functions $\varphi$ on $H$ to genuine
$\tilde\theta$-equivariant invariant functions on $\widetilde H$ by
$(L_D\varphi)(z\,t^*)=b_D\,\tilde\theta(z)\,\Delta(t)\Delta(t^n)^{-1}\varphi(t)$
at good regular elements, and $0$ at the others; it is unitary
(Part~III, Definition~4.1 and Theorem~4.3).
\item[$b_D$, $c_D$] $b_D=|n|_F^{-(r-1)/2}$ and
$c_D=|n|_F^{-(r^2-1)/2}$ (Part~III, \S\S4--5).  Both equal $1$ at a
tame place; in the running example $b_D=3$ and $c_D=27$ at
$\mathfrak p_3$.
\item[Integrality of $L_D$] The question whether $L_D$ carries
irreducible characters to irreducible genuine characters.  It holds at
tame places (Part~III) and is open in general at wild places
(Part~IV; companion paper).
\end{description}

\subsection{Named results and hypotheses}\label{gl:named}
\begin{description}
\item[Theorems A, B, C] Part~I: the comparison of covered elliptic
tori, the matching of regular-supported test functions, and the
global transfer under $(\mathrm H_D)$ (with its converse under (H)).
\item[Theorems D, E, C$'$] Part~II: the equivariant section near the
center and its compatibility with $\jmath$; the matching of arbitrary
genuine test functions; and the global transfer under the weaker
hypothesis $(\mathrm H_D')$ (with its converse under $(\mathrm H')$).
\item[$(\mathrm H_D)$, (H), $(\mathrm H_D')$, $(\mathrm H')$]
Non-cancellation hypotheses: the character of a near-equivalence
packet at the places of $S$ is not identically zero on the relevant
matching test functions (Part~I, Definition~9.9 and
Theorem~9.12; Part~II, Definition~6.4 and \S6.3).  $(\mathrm H_D)$
holds when the $D$-side packet is finite (Part~I,
Proposition~9.10); sufficient conditions for (H) are in Part~I,
Remark~9.13; and $(\mathrm H_D)\Rightarrow(\mathrm H_D')$,
$(\mathrm H)\Rightarrow(\mathrm H')$ (Part~II, \S6.3).
\item[(D1)--(D4)] The steps left open by Part~I: matching of
arbitrary functions, local character relations, finiteness of
near-equivalence packets, and composition with [JC] (Part~I, \S10).
Their status in the later Parts is recorded in Part~I, \S1.3.
\item[Targets A, B, C] The companion paper's three local problems at
the wild place of the running example and beyond (companion paper,
\S1).
\end{description}

%   \subsection{Notation specific to Part I}  ...
%
% The file contains only \subsection's and description lists.  It
% uses only commands from LaTeX, amsmath and amssymb (no macros of
% any Part), so it compiles unchanged in every Part.  Labels carry
% the prefix "gl:" to avoid clashes with the Parts' own labels.
%
% Pointers of the form "Part I, Proposition 4.3" refer to Part I,
% Draft 11 (Part I was reordered in Draft 11; Drafts 8-10 have the
% older numbering, tabulated in Part I's Draft 11 revision log); pointers to Parts II, III
% and IV refer to Drafts 8, 9 and 6.  This file proves nothing.
%
% Version: November 8, 2026 (introduced with Part I, Draft 9; Part I
% pointers renumbered for Part I, Draft 11).
% American English spelling throughout.
% =====================================================================

\subsection{Status labels}\label{gl:labels}
Every statement in the series carries one of the following labels.
A label is never upgraded without a proof, and every change of label
is recorded in the revision log of the draft in which it is made.
\begin{description}
\item[\textbf{[P]}] \emph{Proved here}: a complete argument is
included in the Part in which the label appears, possibly using
results labeled \textbf{[E]}.  A statement proved under an
explicitly stated hypothesis is labeled \textbf{[P]} together with
the hypothesis, as in ``\textbf{[P]} under $(\mathrm H_D)$'', and the
hypothesis itself is listed separately with its own label.
\item[\textbf{[E]}] \emph{Established elsewhere}: a precisely cited
theorem from the literature or from an earlier Part, used as a black
box.  ``\textbf{[P]} (using \textbf{[E]})'' marks a proof that rests
on such inputs; ``\textbf{[E]}/\textbf{[P]}'' marks a statement that
is partly cited and partly proved.
\item[\textbf{[O]}] \emph{Proof outlined}: the steps are identified,
but at least one of them remains to be written out or checked.
\item[\textbf{[C]}] \emph{Conjectural}: new mathematical input is
required.  This includes the goals deferred to later Parts and the
hypotheses under which conditional statements are proved.
\item[\textbf{[P$^{\mathrm c}$]}] \emph{Proved modulo an exhaustive
exact computation}: the proof consists of hand arguments together
with a computation that is exhaustive over an explicitly described
finite set and uses exact arithmetic only (no sampling, and no
floating point without exact re-verification).  Each such script is
committed to the repository, cited by commit hash, and described in
an appendix and in the scripts' README.  At present this label is
used only in the companion paper on the rank-two wild place.
\end{description}

\subsection{Global data}\label{gl:global}
\begin{description}
\item[$K$, $n$, $\mu_n$] $K$ is a totally imaginary number field
containing the group $\mu_n$ of $n$-th roots of unity, and $n\ge1$ is
odd.  All archimedean places of $K$ are complex.
\item[$r$, $D$, $L$, $S$] $r$ is a prime with $(n,r)=1$; $D$ is a
central division algebra over $K$ of dimension $r^2$; $L\subset D$ is
a cyclic maximal subfield, so $[L:K]=r$; $S$ is the (finite) set of
places at which $D$ ramifies.  Since $r$ is prime, every $v\in S$ is
finite and $D_v$ is a division algebra; $|S|\ge2$, and $|S|$ is even
if $r=2$ (Part~I, \S1.1).
\item[$r'$, $c$, $c_v$] Integers with $rr'\equiv1\pmod n$ and
$n\mid r-1+2c$; then $c_v:=r'c$ satisfies $n\mid r-1+2rc_v$ (Part~I,
Assumption~1.1).  Despite the subscript, $c_v$ does not depend on the
place $v$.  The parameter $c$ occurs \emph{before} the relabeling of
$\mu_n$ described in \S\ref{gl:cocycle}, and $c_v$ \emph{after} it.
\item[$\mathbb A_K$, $L_v$] $\mathbb A_K$ is the adele ring of $K$
(written $\mathbb A$ when no confusion is possible), and
$L_v:=L\otimes_KK_v=\prod_{w\mid v}L_w$.
\item[$G_D$] The algebraic group $D^\times$ over $K$, so that
$G_D(K_v)=D_v^\times$.
\item[Running example] $K=\mathbb Q(\zeta_3)$, $n=3$, $r=2$,
$D=(-1,4-\zeta_3)_K$, ramified exactly at
$S=\{\mathfrak p_3,\mathfrak p_7\}$ with $\mathfrak p_3=(1-\zeta_3)$
and $\mathfrak p_7=(3+\zeta_3)$; $L=K(i)=\mathbb Q(\zeta_{12})$,
$c=1$, $r'=2$, $c_v=2$.  The place $\mathfrak p_3$ divides $n$ and
$\mathfrak p_7$ does not.  The note
\texttt{notes/mjl\_running\_example.tex} works the example out across
the series.
\end{description}

\subsection{Local fields and absolute values}\label{gl:local}
\begin{description}
\item[$F$, $\mathcal O_F$, $\varpi$, $q_F$, $p$] $F$ is a
nonarchimedean local field of characteristic zero containing $\mu_n$
(usually $F=K_v$), with ring of integers $\mathcal O_F$, uniformizer
$\varpi$, residue field of order $q_F$, and residual characteristic
$p$.
\item[$|\cdot|_F$, $|\cdot|_p$] $|\cdot|_F$ is the normalized absolute
value, $|\varpi|_F=q_F^{-1}$; on $\mathbb Q_p$ it is written
$|\cdot|_p$, with $|p|_p=p^{-1}$.  In particular $|n|_F=1$ if
$p\nmid n$, and $|F^\times/F^{\times n}|=n^2|n|_F^{-1}$ (Part~I,
proof of Proposition~5.5).
\item[Places dividing $n$; tame and wild places] A place $v$ (or the
field $F=K_v$) \emph{divides $n$} if $p\mid n$.  The places $v\in S$
with $p\nmid n$ are called \emph{tame}, and those with $p\mid n$
\emph{wild}.  In the running example $\mathfrak p_7$ is tame and
$\mathfrak p_3$ is wild ($|3|_{\mathfrak p_3}=3^{-2}$).
\end{description}

\subsection{Hilbert symbols}\label{gl:hilbert}
\begin{description}
\item[$(a,b)_F$] The $n$-th Hilbert symbol of $F$, with values in
$\mu_n$; Part~IV and the companion paper also write $(a,b)_{n,F}$ or
$(a,b)_n$.  It is bimultiplicative, nondegenerate and antisymmetric,
$(a,b)_F=(b,a)_F^{-1}$.  At complex places it is trivial.
\item[Norm compatibility] For a finite extension $F'/F$, $a\in
F^\times$ and $b\in F'^\times$, $(a,b)_{F'}=(a,N_{F'/F}b)_F$; in
particular $(a,b)_{F'}=(a,b)_F^{[F':F]}$ for $a,b\in F^\times$
[Serre, \emph{Local Fields}, Ch.~XIV].
\item[$(a,b)_{\mathbb A}$] For a number field $M\supseteq K$,
$(a,b)_{\mathbb A_M}=\prod_w(a_w,b_w)_{M_w}$.  Hilbert reciprocity:
$(a,b)_{\mathbb A_M}=1$ for $a,b\in M^\times$.
\end{description}

\subsection{Covers and the cocycle normalization}\label{gl:cocycle}
\begin{description}
\item[Cover, $p$, $\overline{H'}$] An \emph{$n$-fold cover} of a
locally compact group $H$ is a topological central extension
$1\to\mu_n\to\overline H\xrightarrow{\ p\ }H\to1$.  For a subgroup
$H'\subset H$ we write $\overline{H'}:=p^{-1}(H')$.  Parts~I and~II
write $\overline H$; Parts~III, IV and the companion paper write
$\widetilde H$, $\widetilde Z$.
\item[Commutators] For lifts $\tilde x,\tilde y$ of commuting
elements, $[x,y]:=\tilde x\tilde y\tilde x^{-1}\tilde y^{-1}\in\mu_n$;
it does not depend on the lifts, and it is bimultiplicative on a
torus [Flicker--Kazhdan, \S3].
\item[Kazhdan--Patterson cocycle $\sigma_F$] The cocycle of
[Kazhdan--Patterson, \S0.1] and [Flicker--Kazhdan, \S2] on
$\mathrm{GL}(r,F)$ with twisting parameter $0$.  On the diagonal
torus,
$\sigma_F(\mathrm{diag}(a_i),\mathrm{diag}(b_i))=\prod_{i<j}(a_i,b_j)_F$.
\item[Parameter-$c'$ cover $\overline G^{(c')}$] The cover of
$\mathrm{GL}(r,F)$ with cocycle
$\sigma_F(g,h)\,(\det g,\det h)_F^{\,c'}$.  The comparison uses
$c'=c_v$; the parameter is forced (Part~I, Remark~5.3).  On the
diagonal torus the commutator of $\mathrm{diag}(t_i)$ and
$\mathrm{diag}(u_j)$ is $\prod_{i,j}(t_i,u_j)_F^{A_{ij}}$ with
$A=-I_r+(1+2c_v)J_r$, where $J_r$ is the all-ones matrix (Part~IV,
\S2.1).
\item[Ambient cover $\overline{\mathcal G}_v$] The cover of
$\mathcal G_v=\mathrm{GL}(r,L_v)=\prod_{w\mid v}\mathrm{GL}(r,L_w)$
with cocycle
$\prod_{w\mid v}\sigma_{L_w}(g_w,h_w)\,(\det g_w,\det h_w)_{L_w}^{\,c_v}$,
the copies of $\mu_n$ being identified (Part~I, \S4.2 and
equation~(1)).  Both covers compared in the series are subgroups of
it, so their common central subgroup $p^{-1}(K_v^\times)$ is literally
the same subgroup on the two sides.
\item[Restriction and relabeling] Pulled back along
$\iota_D:D_v^\times\hookrightarrow\mathcal G_v$, the ambient cocycle
is the local component of the cocycle
$\beta=\sigma_L\cdot(\mathrm{Nrd},\mathrm{Nrd})^c$ of the corrected
1988 paper [JC, equation~(4)].  Pulled back along the diagonal
embedding of $\mathrm{GL}(r,K_v)$, it is isomorphic, \emph{as an
extension}, to the extension with cocycle
$\sigma_F^{\,r}(\det,\det)_F^{\,c}$ (Part~I, Proposition~4.3).  The
group $\mu_n$ is relabeled \emph{once}, at every place and globally,
by $\zeta\mapsto\zeta^{r'}$; after relabeling the second extension is
$\overline G^{(c_v)}$.  Formulas state whether they hold before or
after relabeling.  The isomorphism is one of extensions, not an
equality of cocycles, and it is unique up to
$\tilde g\mapsto\chi(\det p(\tilde g))\,\tilde g$ with
$\chi:F^\times\to\mu_n$ continuous.  Central characters are therefore
never transported through it (Part~I, Remark~4.5).
\item[Scalar centrality] Because $n\mid r-1+2rc_v$, the whole inverse
image $\overline Z=p^{-1}(F^\times)$ of the scalars is central
(Part~I, Lemma~4.2).  A cover with this property is called
\emph{scalar-central}.
\item[Global normalization] Globally one uses the normalized cocycle
$\tau_w(g,h)=\sigma_w(g,h)\,\kappa_w(g)\kappa_w(h)\kappa_w(gh)^{-1}$
of [Flicker--Kazhdan, \S2], where $\kappa_w$ is trivial for $w\mid
n\infty$.  The compact splitting is $\kappa^*_w(k)=(k,1)_\tau$ on
$\mathrm{GL}(r,\mathcal O_{L_w})$ for $w\nmid n\infty$, and the
rational splitting is $\iota(x)=s_\tau(x)\,\kappa_L(x)^{-1}$; in the
unnormalized models the rational splitting is $x\mapsto(x,1)_\sigma$
at every place (Part~I, Lemma~7.1).  The global ambient cover
$\overline{\mathcal G}(\mathbb A_L)$ has cocycle
$\tau_L\cdot(\det,\det)_{\mathbb A_L}^{\,c_v}$.
\item[Star map] For $h$ in the base group and any lift $\tilde h$,
$h^*:=\tilde h^{\,n}$.  It does not depend on the lift, it is
conjugation equivariant, $(ah)^*=a^*h^*$ for scalars $a$, and it is
multiplicative on every torus because $n$ is odd (Part~III,
Definition~2.1; Part~IV, \S2.2).  For $\zeta\in\mu_n(F)$ one has
$\zeta^*=1$.
\end{description}

\subsection{Genuine objects, central characters and good elements}\label{gl:genuine}
\begin{description}
\item[Genuine representation] A representation of $\overline H$ on
which $\mu_n$ acts by the tautological character $\zeta\mapsto\zeta$.
\item[$\tilde\theta$, $\theta$] A genuine character $\tilde\theta$ of
$\overline Z$ (unitary from Part~I, \S7.1 on; globally, a character of
$\overline Z(\mathbb A_K)$ trivial on $\iota(K^\times)$).  The same
$\tilde\theta$ serves both covers, since $\overline Z$ is a common
subgroup.  The ordinary character attached to it is
$\theta(a):=\tilde\theta(a^*)$; it is trivial on $\mu_n$, and
$\theta$ does not determine $\tilde\theta$ (Part~IV, \S2.2).
Part~III requires only that $\tilde\theta|_{\mu_n}$ be faithful.
\item[Genuine test functions] $C_c^\infty(\overline H,\tilde\theta)_{\mathrm{gen}}$
is the space of smooth functions with support compact modulo
$\overline Z$ and $f(z\tilde h)=\tilde\theta(z)^{-1}f(\tilde h)$ for
$z\in\overline Z$; in particular $f(\zeta\tilde h)=\zeta^{-1}f(\tilde h)$.
Characters and invariant class functions of genuine representations
transform by $\tilde\theta$ itself (Part~III, \S2).
\item[Regular, elliptic] An element of $\mathrm{GL}(r,F)$ is
\emph{regular} if its eigenvalues are distinct, and \emph{elliptic
regular} if moreover its characteristic polynomial is irreducible.
Since $r$ is prime, the regular elements of $D_v^\times$ are the
noncentral ones, and they are all elliptic.  An element of a cover is
regular (elliptic) if its image is.
\item[Good element] A regular $\tilde\gamma$ is \emph{good} if it
commutes with $p^{-1}(T_\gamma)$, where $T_\gamma$ is the centralizer
of $\gamma=p(\tilde\gamma)$; equivalently, $\tilde\gamma$ is central
in the covered torus $\overline{T_\gamma}$.  Genuine invariant
functions vanish at regular elements that are not good.  For a
maximal subfield $E$ of $D_v$ (or an elliptic torus $E^\times$ of
$\mathrm{GL}(r,F)$), the good elements of $\overline{E^\times}$ are
$p^{-1}(F^\times E^{\times n})=\overline Z\cdot\{t^*\}$, of index
$|n|_F^{-(r-1)}$, and every genuine irreducible representation of
$\overline{E^\times}$ has dimension $|n|_F^{-(r-1)/2}$ (Part~I,
Proposition~5.5).  Covered elliptic tori are abelian exactly at the
places not dividing $n$.
\end{description}

\subsection{Orbital integrals and matching}\label{gl:matching}
\begin{description}
\item[$\Delta$] For regular semisimple $\gamma$,
$\Delta(\gamma)=|D(\gamma)|_F^{1/2}|\det\gamma|_F^{-(r-1)/2}$, where
$D(\gamma)$ is the discriminant of the (reduced) characteristic
polynomial and $\det=\mathrm{Nrd}$ on the division-algebra side.
Thus $\Delta(\gamma)^2$ is the Weyl discriminant
$|\det(1-\mathrm{Ad}(\gamma^{-1}))|_{\mathfrak g/\mathfrak t}|_F$
(Part~I, Definition~6.1 and Lemma~6.12).
\item[$\Phi$, $F=\Delta\Phi$] $\Phi(\tilde\gamma,\tilde f)$ is the
orbital integral over $T_\gamma\backslash G$ of a genuine test
function at a good regular $\tilde\gamma$, and
$F(\tilde\gamma,\tilde f)=\Delta(\gamma)\Phi(\tilde\gamma,\tilde f)$
is the normalized orbital integral, set to $0$ at regular elements
that are not good (Part~I, Definition~6.1).  Part~III writes
$J(t,f)=\Delta(t)\,O(t,f)$ for the same function, with its own
measure conventions.
\item[Canonical identification $\jmath$] At $v\in S$ and for a maximal
subfield $E\subset D_v$ with an embedding $j:E\hookrightarrow
M_r(F)$, every isomorphism of covered tori that is the identity on
$\overline Z$ has the same restriction $\jmath_{E,j}$ to the good
elements; these restrictions induce a bijection $\jmath$ from the good
regular classes of the cover of $D_v^\times$ to the good regular
elliptic classes of the cover of $\mathrm{GL}(r,F)$ (Part~I,
Theorem~5.9 and Corollary~5.12).
\item[Matching] Genuine test functions $\tilde f_D$ and $\tilde f$
\emph{match} if (M1)
$F(\tilde\gamma,\tilde f_D)=(-1)^{r-1}F(\jmath(\tilde\gamma),\tilde f)$
at every good regular $\tilde\gamma$, and (M2) $F(\tilde\delta,\tilde
f)=0$ at every good regular non-elliptic $\tilde\delta$ (Part~I,
Definition~6.3).  The sign sits in the matching of functions, so that
the expected character relation is
$\Theta_{\overline\pi_D}=(-1)^{r-1}\Theta_{\overline\pi}\circ\jmath$.
\item[Supercusp form] A genuine test function on the cover of
$\mathrm{GL}(r,F)$ whose integrals along the unipotent radicals of
all proper parabolic subgroups vanish, for example a matrix
coefficient of a supercuspidal representation [Flicker--Kazhdan,
\S13].
\item[Germ] The asymptotic behavior of orbital integrals near a
singular element, expressed through Shalika germs [Flicker--Kazhdan,
\S8]; used near the center in Part~II.
\end{description}

\subsection{Haar measures}\label{gl:measures}
\begin{description}
\item[Part~I] At $v\in S$ the Haar measures on $D_v^\times/F^\times$
and $\mathrm{GL}(r,F)/F^\times$ are arbitrary but fixed, and matching
is defined with respect to them; torus measures are transported by
conjugation and do not affect matching.  At $v\notin S$ the measures
correspond under the identifications $\psi_v$ (Part~I, \S\S6.1, 7.3).
\item[Part~II] At $v\in S$,
$\mathrm{vol}(D_v^\times/F^\times)=d(\mathrm{St})^{-1}$, where
$d(\mathrm{St})$ is the formal degree of the Steinberg representation
for the chosen measure on $\mathrm{GL}(r,F)/F^\times$ (Part~II,
Definition~2.1).
\item[Part~III] For the transform $L_D$, the quotients $H/Z$ and
$T/Z$ carry their Haar \emph{probability} measures (Part~III, \S4).
\item[Formal degree] For a representation $\sigma$ that is square
integrable modulo the center, the constant $d(\sigma)$ in the Schur
orthogonality relations; it depends on the Haar measure.
\end{description}

\subsection{The division-side transform}\label{gl:transform}
\begin{description}
\item[$H$, $\widetilde H$, $L_D$] For a local division algebra $D$ of
degree $r$ over $F$, $H=D^\times$ and $\widetilde H$ is its cover from
Part~I.  The transform $L_D$ sends $\theta$-equivariant invariant
$L^2$ functions $\varphi$ on $H$ to genuine
$\tilde\theta$-equivariant invariant functions on $\widetilde H$ by
$(L_D\varphi)(z\,t^*)=b_D\,\tilde\theta(z)\,\Delta(t)\Delta(t^n)^{-1}\varphi(t)$
at good regular elements, and $0$ at the others; it is unitary
(Part~III, Definition~4.1 and Theorem~4.3).
\item[$b_D$, $c_D$] $b_D=|n|_F^{-(r-1)/2}$ and
$c_D=|n|_F^{-(r^2-1)/2}$ (Part~III, \S\S4--5).  Both equal $1$ at a
tame place; in the running example $b_D=3$ and $c_D=27$ at
$\mathfrak p_3$.
\item[Integrality of $L_D$] The question whether $L_D$ carries
irreducible characters to irreducible genuine characters.  It holds at
tame places (Part~III) and is open in general at wild places
(Part~IV; companion paper).
\end{description}

\subsection{Named results and hypotheses}\label{gl:named}
\begin{description}
\item[Theorems A, B, C] Part~I: the comparison of covered elliptic
tori, the matching of regular-supported test functions, and the
global transfer under $(\mathrm H_D)$ (with its converse under (H)).
\item[Theorems D, E, C$'$] Part~II: the equivariant section near the
center and its compatibility with $\jmath$; the matching of arbitrary
genuine test functions; and the global transfer under the weaker
hypothesis $(\mathrm H_D')$ (with its converse under $(\mathrm H')$).
\item[$(\mathrm H_D)$, (H), $(\mathrm H_D')$, $(\mathrm H')$]
Non-cancellation hypotheses: the character of a near-equivalence
packet at the places of $S$ is not identically zero on the relevant
matching test functions (Part~I, Definition~9.9 and
Theorem~9.12; Part~II, Definition~6.4 and \S6.3).  $(\mathrm H_D)$
holds when the $D$-side packet is finite (Part~I,
Proposition~9.10); sufficient conditions for (H) are in Part~I,
Remark~9.13; and $(\mathrm H_D)\Rightarrow(\mathrm H_D')$,
$(\mathrm H)\Rightarrow(\mathrm H')$ (Part~II, \S6.3).
\item[(D1)--(D4)] The steps left open by Part~I: matching of
arbitrary functions, local character relations, finiteness of
near-equivalence packets, and composition with [JC] (Part~I, \S10).
Their status in the later Parts is recorded in Part~I, \S1.3.
\item[Targets A, B, C] The companion paper's three local problems at
the wild place of the running example and beyond (companion paper,
\S1).
\end{description}
|, and \S10.11 ``Notation specific to Part~I''
(about twenty entries with pointers).  The entries of the old informal
glossary are all kept: series-wide ones (cover, genuine, ambient
cover, good element, elliptic regular element, canonical
identification, normalized orbital integral, matching, supercusp form,
formal degree, germ) in the shared file, and Part-specific ones
(near-equivalence packet, corner Hecke algebra, elliptic
detectability) in \S10.11.  Part~I must now be compiled from
\texttt{part-I/} so that \texttt{../glossary.tex} is found; the
repository README was updated accordingly.

\item \textbf{(b) Summary of results.}  New one-page unnumbered section
after \S1 (it appears in the table of contents).  Seven items with
labels and pointers: the ambient cover; Theorem~A; Theorem~B and the
failure near the center; rational lifts and the trace formulas; the
packet identity; Theorem~C, its converse and Lemma~7.8; and (D1)--(D4)
with their status in Parts~II--III; plus one line on the running
example.  The labels are those already in the paper.

\item \textbf{(c) Literature review.}  \S11 was rewritten as a review
for the whole series, in five paragraphs: the Jacquet--Langlands
correspondence; covering groups; metaplectic correspondences; central
extensions of division algebras; and the present series.  Parts~II--IV
will cite it and add only what is specific to them.  The claims of the
old \S11 are kept (in particular the zbMATH search over the papers
citing [FK]).  Twenty-six references were added: BD, FK88, Fli80,
Gel76, GG18, GGW, JL, Kaz84, Klo86a, Klo86b, Kub67, Kub69, Li1, Li2,
Mez01, Mez02, Mez04, PRa88, PRp96, Rap88, Rog83, Sav88, Shi73, Wal80,
Wal91, Wei18.  The new text cites no theorem numbers from them.

\emph{Reference checks (bibliographic data only).}
\begin{itemize}
\item Checked this session against numdam/EuDML/publisher pages:
Moore (Publ.\ IH\'ES 35 (1968), 5--70; Part~I's entry is
\emph{correct}), Prasad--Rapinchuk (84 (1996), 91--187), Li~II (Ann.\
ENS 45 (2012), 787--859), Flicker--Kazhdan 1988 (J.\ Analyse 50,
189--200), Mezo~I (CJM 54 (2002), 92--137), Mezo~II (Represent.\
Theory 5 (2001), 524--580), Shimura (Ann.\ 97 (1973), 440--481),
Gelbart (LNM 530, 1976), Gan--Gao and Gan--Gao--Weissman
(Ast\'erisque 398, 187--275 and 1--31), Klose 1986a (Math.\ Z.\ 191,
261--282).
\item Checked only through a citing bibliography (the reference lists of
Prasad--Rapinchuk, Mezo or Moore on numdam and the publishers' pages): Prasad--Raghunathan 1988,
Rapinchuk 1988, Kazhdan 1984, Flicker 1980, Kubota 1967.
\item Klose 1986b: numdam/geodesic gives pages 625--650,
Prasad--Rapinchuk's list gives 625--649; we used 625--650.
\item Mezo 2004 (Pacific J.\ Math.\ 215, no.~1): the first page is 85,
but we could not confirm the last page, so the entry gives no page
range.
\item Li~I: we could not confirm the journal; the entry cites
arXiv:1004.4011.
\item Weissman (Ast\'erisque 398): first page 33 from the volume's
table of contents; last page 186 inferred from Gan--Gao starting at
187.
\item Taken from the series' existing bibliographies without a new
check: Brylinski--Deligne, Jacquet--Langlands, Kubota 1969, Rogawski,
Savin, Waldspurger 1980 and 1991.
\item Contents: the descriptions of the new references are at the level
of their titles, abstracts and front matter.  The originals of
Klose's papers, Savin, and Mezo~II were not available to us.
\end{itemize}

\item \textbf{(d) Lead-ins.}  Every section (\S\S1--9 and
Appendices~A, B) now says, in two to four sentences, what it proves,
what it depends on and what notation it introduces; \S\S10 and~11
have new lead-ins.  Subsection lead-ins were extended where the
subsection introduces notation or a hypothesis used later
(\S\S2.2, 3.4, 3.5, 4.1, 4.3, 5.1, 5.3, 6.3, 7.1).  The existing
one-sentence lead-ins of the other subsections were kept as they were.

\item \textbf{Introduction.}  The \S1 lead-in points to the Summary and
to \S10; \S1.4 (Conventions) and \S1.5 (labels) point to the matching
subsections of \S10; \S1.6 (Organization) describes the new
\S\S10--11 and the Summary.  The text of \S\S1.4--1.5 is otherwise
unchanged.

\item Title line: Draft~9, November 8, 2026.  Authors and acknowledgments
unchanged.
\end{enumerate}

\section*{Points for David (not acted on)}
\begin{enumerate}[leftmargin=*]
\item \textbf{Meaning of [O].}  The request describes \textbf{[O]} as
``open''.  Parts~I--IV, the companion paper and the README all define
it as ``proof outlined but not fully established'', and in Part~I open
goals are labeled \textbf{[C]}.  The glossary keeps the papers' wording;
changing it would relabel statements, so we have not done so.
\item \textbf{Flicker--Kazhdan, ``A simple trace formula'' (1988).}
Its abstract says it extends the metaplectic correspondence, as well as
the simple-algebras correspondence, to automorphic forms with a single
supercuspidal component.  This may overlap Proposition~6.2 (our
variant of [FK, \S18] with one supercusp form).  Whether to cite it as
an alternative source, or to compare the two in a remark, is a
mathematical decision; the literature review cites it for its approach
only.
\item \textbf{Klose's papers.}  Klose (a student of Patterson) studied
metaplectic extensions of quaternion division algebras (Math.\ Z.\ 191
and 193, 1986).  This is prior work on covers of $D^\times$.  The
companion paper's literature section says that we are ``not aware of
an independent literature on $n$-fold covers of division algebras
beyond [JC] and this series''; that sentence should be revised when the
companion paper is next edited.  Whether Klose's covers agree with
those of [JC] for $r=2$ has not been checked.
\item \textbf{Companion paper bibliography.}  Its entry for Moore gives
pages 157--222; the correct range is 5--70.
\item \textbf{Subsection lead-ins.}  The request asks for subsection
lead-ins only where notation or a hypothesis is introduced.  Part~I
already had a one-sentence lead-in for every subsection (an earlier
request); we kept them.  They can be pruned if you prefer.
\item \textbf{Placement.}  The Summary is unnumbered and the glossary
sits at \S10 (not at the front) only to keep Part~I's numbering fixed
for Parts~II--IV.  If you would rather have the glossary at the front,
the cross-references in Parts~II--IV would have to be updated in the
same pass.
\end{enumerate}

\section*{Files}
\texttt{glossary.tex} (new, repository root);
\texttt{part-I/mjl\_part1\_draft9.tex} and \texttt{.pdf};
this log (\texttt{part-I/mjl\_part1\_revision\_log\_draft9.tex} and
\texttt{.pdf}); \texttt{README.md} (Part~I draft number, build note,
labels table).  Draft~8 is kept unchanged.  No scripts were added or
changed.
\end{document}
